\scanpage{8}

5) Soit $K$ un fermé de l'\unsure{espace} $X$, $(\deleted{\unsure{K}}\unsure{U}{}^{\deleted{\unsure{\prime}}}_i{}_{\deleted{\unsure{k|K}}})_{i\in I}$ une famille d'ouverts relatifs de $K$, mutue\uline{llement disj}oints, tels que $\dot K\subset\bigcup_i U_i{}_{\deleted{k}}$. \inserted{Soit $B_i=\dot K\cap U_i$, de sorte que $\dot K=\coprod_i B_i$.} \deleted{Supposons \unsure{q}} Pour $\forall i$, soit $C_i=\overline U_i{}_{\deleted{k}}\smallsetminus U_i{}_{\deleted{k}}$ la frontière relative de $U_i{}_{\deleted{k}}$ dans $K$, supposons que $C_i\subset\bigl(K-\overline U_i{}_{\deleted{k}}\bigr)^{\! -}$ (cf\ldots). Posons $\deleted{\unsure{X}}\deleted{\unsure{U_i}}X_i=X\smallsetminus U'_i$. Alors
\[
\left\{\begin{aligned}
&X_i\ \text{est loc fermé}\\
&\operatorname{Int}X_i=X\smallsetminus\overline{U'_i}=X_i\smallsetminus C_i\\
&\overline X_i=X_i\cup \unsure{\dot B'_i|B'_i}\qquad(B'_i=\dot K\smallsetminus B_i)
\end{aligned}\right.
\]
% slanted
\margin{\deleted{\unsure{i.e.}\ $C_i$}}

$\operatorname{Int}X_i$ et \uline{$U_i\smallsetminus \unsure{C_i|C'_i}$} sont \deleted{deux} \unsure{ouverts} de $X$ dont chacun est le \unsure{complémentaire} de l'adhérence de l'autre.

Soit $W=X\smallsetminus K$, \deleted{\unsure{semble}} donc \inserted{$\dot W={}$} $\overline W\cap K=\dot K=\coprod B_i$, l'inclusion $W\subset X_i$ est une \unsure{imm} \unsure{ouvert|ouverte} et $\dot W\cap X_i=\dot K\cap X_i=B_i$; \deleted{$\unsure{X_i}$} i.e \deleted{\unsure{l'on}} l'adh. de $W$ dans $X_i$ est $W\cup B_i$, donc la \deleted{bord}\inserted{frontière} de $W$ dans $X_i$ est $B_i$. On a
\[
X_i=(\overline W\cap X_i=W\cup B_i)\amalg_{B_i}\overline U_i
\]

% margin note written sideways (upside down)
\margin{\deleted{\unsure{Monsieur} le \unsure{P|D}}}

Alors \unsure{la|le} somme \unsure{amalgamée} de tous les \unsure{$X_i$} sur $W$ est un \unsure{espace}
% Last line meets the lower edge of the scan.
% archive sheet number: 7
